% Generated draft: requires mathematical and source review.
\documentclass[11pt]{article}
\usepackage[margin=25mm]{geometry}
\usepackage{fontspec}
\setmainfont{Latin Modern Roman}
\usepackage{amsmath,amssymb,mathtools,unicode-math}
\setmathfont{Latin Modern Math}
\usepackage{enumitem,needspace,xurl,hyperref}
\setlength{\parindent}{0pt}
\setlength{\parskip}{5pt}
\setlength{\emergencystretch}{3em}
\newcommand{\N}{\mathbb{N}}
\newcommand{\Z}{\mathbb{Z}}
\newcommand{\Q}{\mathbb{Q}}
\newcommand{\R}{\mathbb{R}}
\newcommand{\C}{\mathbb{C}}
\newcommand{\F}{\mathbb{F}}
\newcommand{\A}{\mathbb{A}}
\newcommand{\PP}{\mathbb{P}}
\newcommand{\E}{\mathbb{E}}
\newcommand{\eps}{\varepsilon}
\newcommand{\dd}{\,\mathrm{d}}
\begin{document}
\section*{UnsolvedMath ID: 30002881}
\textbf{Source problem code:} catalogue entry
\subsection{Definitions and mathematical statement}
Let \(p\) be a prime, let \(G\) be a finite group, and let \(S\) be a Sylow \(p\)-subgroup of \(G\). The supplied statement does not specify the field \(k\), its characteristic, or whether \(S\ne 1\). The following is a conditional formulation of the intended modular representation setting: assume that \(k\) is a field of characteristic \(p\) and that \(S\ne 1\). These additional hypotheses require confirmation; any intended splitting-field or algebraic-closure hypothesis is also unresolved.

For a subgroup \(Q\le G\), write
\[
N_G(Q)=\{g\in G:gQg^{-1}=Q\},
\]
and write \(H'=[H,H]\) for the commutator subgroup of a group \(H\). Interpreting the source's subgroup indices as nontrivial subgroups of \(S\), define, for every \(1\ne Q\le S\),
\[
\rho_1(Q)=N_G(Q)',\qquad
\rho_i(Q)=\left\langle\bigcup_{1\ne R\le S}\bigl(N_G(Q)\cap\rho_{i-1}(R)\bigr)\right\rangle
\quad(i\ge2),
\qquad
\rho^\infty(S)=\bigcup_{i\ge1}\rho_i(S).
\]
Here \(\langle X\rangle\) denotes the subgroup generated by \(X\). The displayed initial term and recursion are retained exactly as supplied; the source does not expressly give the domain of \(Q\) or a convention for \(S=1\).

All modules below are finite-dimensional left modules. A \(kH\)-module \(M\) is endotrivial if
\[
\operatorname{End}_k(M)\cong k\oplus P
\]
as \(kH\)-modules for some projective \(kH\)-module \(P\), where \(k\) is the trivial module and \(H\) acts on endomorphisms by \((h\cdot f)(m)=h f(h^{-1}m)\). Write \(U\simeq_{\mathrm{st}}V\) if \(U\oplus P\cong V\oplus P'\) for some projective modules \(P,P'\) over the same group algebra.

\textbf{Supplied conjecture, subject to the qualifications above.} For every one-dimensional \(kN_G(S)\)-module \(L\), with associated character \(\lambda_L:N_G(S)\to k^\times\),
\[
\left(\exists\text{ an endotrivial }kG\text{-module }M:
\operatorname{Res}^{G}_{N_G(S)}M\simeq_{\mathrm{st}}L\right)
\quad\Longleftrightarrow\quad
\rho^\infty(S)\le\ker(\lambda_L).
\]

The source also calls this equivalent to saying that the image of \(K(G)\to K(N_G(S))\) is the dual of \(N_G(S)/\rho^\infty(S)\), but it defines neither \(K\) nor the dual. Under the conventional interpretation, let \(T(H)\) be the group of stable isomorphism classes of endotrivial \(kH\)-modules, with tensor product as its operation, and set
\[
K(H)=\ker\!\left(\operatorname{Res}^{H}_{S}:T(H)\longrightarrow T(S)\right)
\quad\text{for }H\in\{G,N_G(S)\}.
\]
The intended dual would then be
\[
\operatorname{Hom}\!\left(N_G(S)/\rho^\infty(S),k^\times\right),
\]
viewed through the stable classes of the corresponding one-dimensional modules. With these interpretations the source's second formulation reads
\[
\operatorname{im}\!\left(\operatorname{Res}^{G}_{N_G(S)}:K(G)\longrightarrow K(N_G(S))\right)
=
\left\{[L]\in K(N_G(S)):
\dim_k L=1,\ \rho^\infty(S)\le\ker(\lambda_L)\right\}.
\]
The conventions for \(K\), the dual, and the asserted equivalence require confirmation from the cited source.
\subsection{Short English statement}
The supplied conjecture asks whether a one-dimensional representation of the normalizer of a Sylow \(p\)-subgroup extends stably to an endotrivial representation of the whole finite group exactly when it kills the recursively defined subgroup \(\rho^\infty(S)\). The field hypotheses, the trivial Sylow case, and the notation in the claimed equivalent formulation need confirmation. The supplied curation labels the problem partially solved; that assessment is not independently verified here.
\subsection{Sources}
\begin{itemize}
\item Local UnsolvedMath registry, \texttt{data/unsolvedmath/problems.json}: supplied complete record with permanent ID \(30002881\), title \emph{Restriction Criteria for Endotrivial Modules}, and exact source code \texttt{catalogue entry}. The record gives proposed year \(2015\) and identifies no individual proposer.
\item Supplied report citation: \emph{Cohomology of Finite Groups: Interactions and Applications} (2016), Oberwolfach Reports, pp.\ 1317--1384, DOI \texttt{10.4171/owr/2015/24}. \url{https://doi.org/10.4171/owr/2015/24}. The supplied literature list also identifies this as \emph{Oberwolfach Report 2015/24}.
\item Supplied literature reference: Carlson--Th\'evenaz, \emph{The torsion subgroup of endotrivial modules}. \url{https://arxiv.org/abs/1408.3982}.
\item Literature assessment dated exactly \texttt{2026-08-21}: the supplied curation labels the problem partially solved and reports that the conjecture is proved for groups with abelian Sylow \(p\)-subgroups, that further special Sylow cases are known, and that no general proof is established. This is supplied dated curation derived from the OWR record, not an independently verified literature assessment; the record itself requests independent verification of time-sensitive status and cited claims.
\item The registry labels the statement \texttt{corrected\_verified} and says that mathematical content was preserved while notation, typography, or omitted definitions were repaired for publication. This is a supplied editorial label, not mathematical verification of this draft or independent inspection of either reference.
\end{itemize}
\end{document}
